本次软件开发实践中，我以前端初学者的身份参与 WEMOVE SPORTS 网站重构，负责站点公共界面、%
文章与品牌内容、常见问题、媒体和资料下载等工作。面对页面布局、交互效果和内容管理任务，我首先调研
HTML、CSS、JavaScript 及其与前端框架、工程工具之间的关系，再结合项目理解 Vue 3、TypeScript、Vue
Router、Pinia、Vite 和浏览器 API 的具体作用。在完成基本页面的基础上，%
我进一步围绕拍立得商品卡片的三维展示、文章与图片的加载、图文内容的编辑三个问题展开技术调查。

报告按照“认识技术—结合任务使用—向相关技术扩展”的顺序，比较 CSS 三维变换、Three.js、WebGL 与 WebGPU
的适用情况，分析客户端渲染、服务端渲染、静态生成、懒加载和响应式图片的优缺点，%
并讨论从原生可编辑区域到结构化编辑器、协同编辑的发展。通过调研，我逐渐认识到，%
前端开发不仅是把页面做得美观，还要协调内容结构、交互状态、设备能力和维护成本。%
通过本项目的开发，我对适度使用三维效果、分清不同资源的加载时机、约束图文结构有了更深入的理解。
\vspace{1.5cm}

\noindent\textbf{关键词：}前端入门；Vue 3；三维展示；内容加载；富文本编辑；Web 开发


